\chapter{Theoretical Framework}

This chapter outlines the foundations supporting the development of
\emph{SAKLOLO: A Hybrid Semantic-Probabilistic Framework for Analysing Social
	Media Emergency Feeds in Bulacan and CAMANAVA}. It is grounded in Vector Space
Semantics, the Distributional Hypothesis, and Bayesian Generative Theory,
justifying the probabilistic extraction of semantic structures from
unstructured disaster text. The literature review discusses Count
Vectorization mechanics, contextual language representations (\xlmr),
generative large language models (LLMs), and Latent Dirichlet Allocation (LDA)
applications in crisis environments. A synthesis highlights gaps in processing
mixed-language (Taglish) emergency feeds and guides the proposed mathematical
joint matrix model.

\section{Related Theories}

The conceptual foundation of this research is grounded in Sensemaking Theory,
which asserts that in digital crisis environments, offloading the initial
chaotic, high-velocity data to automated systems allows human responders to
bypass information overload and focus directly on high-level resource
allocation \parencite{pezanowski2022}. To mathematically achieve this automated
sensemaking, the computational framework relies on the Distributional Hypothesis
in computational semantics. \parencite{balkir2016} explain that this hypothesis
assumes words occurring in similar contexts inherently share similar semantic
meanings. This provides the theoretical justification for mapping text into
vector spaces. It theorizes that even highly disorganized human crisis
communication contains latent geometric patterns, allowing related disaster
concepts to naturally cluster together when analyzed algebraically.

Complementing this geometric assumption is the statistical principle of
Exchangeability, rooted in De Finetti's theorem. \parencite{blei2003} utilized
this theory to establish that when modeling text data, the joint probability of
a sequence of terms is invariant to their specific grammatical order. This
provides the core theoretical justification for the ``bag-of-words'' approach
used in probabilistic topic modeling. Because the theory assumes that documents
are the end result of a probabilistic generative process rather than rigid
grammatical structures, the proposed system is mathematically validated in
ignoring the broken syntax and informal structures typical of emergency posts,
focusing purely on the statistical distribution of variables to uncover hidden
situational categories.

\section{Review of Related Literature}

\subsection{Computational Frameworks in International Disaster Literature}

The theoretical foundational framework for processing unstructured crisis
communications relies on computational linear algebra, vector space modeling,
and probabilistic generative models. According to \parencite{mihunov2022}, vector
space mechanics represent qualitative narrative posts as numeric constructs in
high-dimensional matrices to enable systematic computation during disasters.
Within this space, a document--term matrix captures documents across unique
vocabulary tokens, where raw text is mathematically transformed using Count
Vectorization to systematically extract critical crisis keywords for downstream
processing.

Complementing these deterministic vector decompositions, probabilistic
frameworks model the underlying semantic structure of text.
\parencite{blei2003} introduced Latent Dirichlet Allocation (LDA) as a Bayesian
generative probabilistic model that views each document as a multinomial
distribution over latent topics, which are themselves multinomial
distributions over vocabulary words governed by Dirichlet priors. Expanding on
this by bridging probabilistic topic models with deep learning,
\parencite{zang2026} highlights the integration of LDA with contextual transformer
architectures (such as BERT). Inferring hidden topic assignment variables
alongside deep contextual embeddings from observed social text feeds allows for
highly accurate, granular identification and tracking of emergent situational
themes as rapid disaster events unfold.

\subsection{Social Media Utilization in Philippine Disaster Risk Reduction}

A literature review on the role of social media in disaster response in the
Philippines found that it was heavily employed in response to typhoons and
earthquakes, but other disasters such as COVID-19, floods, and vegetation
fires are not fully explored. The review indicated that the usage of social
media was heavy during response and recovery but minimal during preparation.
Because misinformation, noise, and infrastructure limitations remain a concern
for its utilization as a reliable disaster response tool, a hybrid
computational approach is recommended. Social media was found to be highly
effective in coordinating recovery and updates, provided the data is
computationally filtered to separate the informative from the controversial
\parencite{inferio2025}.

The Philippines has four thematic areas of Risk Reduction and Management.
Disaster Prevention and Mitigation is aimed at strengthening community
resilience by minimizing exposure to risks and increasing recovery capacity.
Disaster Preparedness involves strengthening community and institutional
capacity via collaborative policies and systems. Disaster Response focuses on
the reduction of preventable casualties and ensures the basic needs of
communities are met and restored rapidly. Finally, Disaster Rehabilitation and
Recovery deals with the restoration of livelihoods, economic opportunities, and
essential infrastructure that supports the physical and psychological
rehabilitation of those affected \parencite{inferio2025}.

The North Manila Bay Delta consists of the downstream part of the rivers of the Pampanga River Basin, including the coastal area. The wider basin has a catchment area of $10,000\text{ km}^2$ and is roughly divided into three river systems. These are the Pampanga main river system (about $8,000\text{ km}^2$), the Angat River system (about $1,085\text{ km}^2$), and the Pasac river system (about $1,371\text{ km}^2$), which originates from Mt. Pinatubo and is connected to the Pampanga River by the Bebe-San Esteban Cut-off Channel \parencite{RoyalHaskoningDHV}.

This delta covers almost $3,000\text{ km}^2$ of the southern central valley of Luzon, stretching from Pampanga to Bataan, Bulacan, and the CAMANAVA area of Metro Manila. The plains are so flat and close to sea level that the 1-meter elevation line extends 10 to $20\text{ km}$ inland, so even a small rise in relative sea level translates to a large encroachment of the sea. The region is also a major rice, sugar, and fish producer, and groundwater is its main source of fresh water for households, industries, fishponds, and farms \parencite{Rodolfo2003}.

\subsection{Hydrometeorological Characteristics}

The coast is tide-dominated and has a mixed, predominantly diurnal tide with an average range of $1.2\text{ meters}$ during spring tide and $0.4\text{ meters}$ during neap tide. During the southwest monsoon from June to September, the tidal amplitude can increase by 80 percent, and wind waves of over $3\text{ meters}$ may occur. The wet season runs from June to November, with average monthly rainfall of 180 to $410\text{ mm}$. The area is affected on average by one to two cyclone-related storm surges per year, and rainfall at the peak of the cyclone period can reach $400\text{ mm}$ a day. The relative sea-level rise in Manila Bay is on average 1.5 to $2.5\text{ cm}$ per year, and land subsidence caused by groundwater over-extraction is occurring in several areas \parencite{RoyalHaskoningDHV}.

\subsection{Linguistic Context and Geographic Scope Selection}

There are four main Tagalog dialects: Northern Tagalog Dialects of Bulacan (Bulacan \& Nueva Ecija) and Bataan (Bataan \& Zambales), Central Tagalog Dialects of Laguna (Laguna \& Rizal) and Metro Manila, Southern Tagalog Dialects of Batangas (Cavite \& Batangas), Lubang (Mindoro), and Tayabas (Aurora, Camarines Norte, Camarines Sur, \& Quezon), and lastly Tagalog Dialects of Marinduque \parencite{Adriano2022}.

The North Manila Bay Delta shares common hydrometeorological traits of natural disaster frequencies and structural instances. It consists of minor parts of Bataan, Pampanga, Bulacan, and CAMANAVA. Bataan was omitted from this study due to minor inclusions on the delta; while Pampanga was omitted due to not being included in major Tagalog dialects and having a different linguistic profile among its speakers. That leaves Bulacan and CAMANAVA as the geographic scope of the study due to sharing the same hydrometeorological profile and relatively similar linguistic profile of Tagalog (Northern and Central dialect), English, and Taglish.

\section{Review of Related Studies}

\subsection{Global Evaluations of Text Processing in Disaster Management}

Empirical implementations across Asian and global disaster contexts demonstrate
the practicality of combining probabilistic topic modeling with
transformer-based deep learning to process high-velocity unstructured social
text feeds. \parencite{zang2026} utilized the LDA topic modeling algorithm
alongside prompt-BERT models to process unorganized short-text feeds during the
2023 Hebei Floods in China, proving that hybrid architectures effectively
categorize heterogeneous social discourse into actionable thematic clusters for
crisis managers.

Evaluating short-text social streams during major hydrometeorological events
further highlights the necessity of strict statistical metric validation.
\parencite{anggara2025} applied LDA topic modeling to Twitter data during the 2024
West Sumatra flash floods, demonstrating that evaluating topic coherence
scores ($C_U$) across various topic ranks ($K$) isolates the optimal matrix
structure, ensuring distinct and semantically meaningful categories without
topic overlap.

In Western disaster management contexts, feature extraction and topic modeling
have proven vital for real-time situational awareness. \parencite{mihunov2022} used
LDA to analyze disaster impact feeds during Hurricane Harvey, successfully
isolating 9 high-priority target topic vectors out of 24 computed themes to
directly map physical infrastructure damage, power grid outages, and localized
civilian entrapment locations.

Further demonstrating the efficacy of localized text extraction during severe
flooding in Southeast Asia, \parencite{rahmadan2020} investigated the Jakarta flash
flood domain by applying Latent Dirichlet Allocation (LDA) to high-velocity
Twitter streams. Their research successfully clustered unstructured text feeds
by localized impact, utilizing the algorithm to map explicit target tokens that
detailed specific flooded neighborhoods, severe water depth thresholds, and
immediate rescue dispatch needs. Because the hydrology and social media usage
in Jakarta closely mirror the urban--rural challenges found in Philippine
municipalities, this study proves that probabilistic generative models can
segment baseline social media noise from urgent environmental disaster events
to directly support local municipal response.

\subsection{Local Applications of Topic Modeling and Sentiment Analysis in
	Crisis Studies}

In a local context, Visayan scholars from the University of San Carlos of Cebu
and Cebu Technological University utilized Latent Dirichlet Allocation (LDA)
and the Hierarchical Dirichlet Process (HDP) to model a qualitative risk
assessment for disaster risk reduction response. They attained qualitative data
via the disaster risk reduction suggestions of Malasakit, a crowd-sourced
disaster initiative across barangays focusing on how to service vulnerable
populations \parencite{gorro2021}.

Before conducting the LDA, they performed data cleaning to remove special
characters, stop words, and strings under four letters. They then vectorized
their entries as a computational prerequisite. For LDA, they measured
perplexity scores to determine how stable the machine learning model was in the
face of uncertainty, where a lower score indicated better performance. After
fitting the models, they discovered the general structural trends: solid waste
management, relief and emergency planning, and early warning systems were the
most heavily weighted topic sets \parencite{gorro2021}. This study provides a
quantitative approach to determining the optimal number of topics based on
perplexity scores, which is a valuable methodological addition to local
studies.

Another local study utilized LDA in the context of the COVID-19 pandemic in
Cagayan Valley, placing a specific emphasis on sentiment analysis. The study
utilized Facebook posts from government agencies between March and June 2020,
applying standard preprocessing protocols. Using a Java package called Mallet,
they fed the data into an LDA model to extract core topics. The generated topic
sets went through labeling to identify common themes with the help of the
Municipal Disaster Risk Reduction and Management Office (MDRRMO) of Saguday,
Quirino, utilizing six annotated tests judged by multiple public and private
sector experts. The study identified five overarching themes capturing public
sentiment: COVID-19 testing, relief operations, contact tracing, government
programs, and community quarantine \parencite{pagudpud2020}.

A separate study conducted by Silliman University researchers in Dumaguete City
sourced qualitative data from Twitter (X) to conduct topic modeling and
document clustering regarding localized sentiments. Utilizing the Twitter4j
API, the program gathered 99,942 tweets between 2015 and 2016. After
preprocessing the mixed-language texts (Tagalog and English) utilizing
string-searching and dictionary-matching algorithms (Aho--Corasick and
Damerau--Levenshtein Distance), 22,633 valid tweets remained for training and
testing \parencite{montenegro2018}. The researchers utilized Gibbs Sampling to
incorporate observed data into the LDA sampling process. Finally, they
conducted sentiment analysis via Support Vector Machines (SVM) to classify the
texts into positive, negative, or neutral sentiments. While the SVM approach
proved functional, it highlights the methodological progression toward modern
hybrid architectures and joint matrices that provide deeper contextual
understanding than classical SVMs \parencite{montenegro2018}.

\subsection{Related Undergraduate Theses in Bulacan State University}

Bulacan State University has been instrumental in the development of
technological applications related to disaster management through various
undergraduate studies. An example of this is the study entitled
\emph{LokaMatika: Two-Way Location Sharing and Automated Emergency Response
	Systems with Trilateration and Kinematics} \parencite{manalo2024}. This study
designed a mobile application utilizing Dart and Flutter that allows two-way
location sharing to help users report and track the locations of vulnerable
individuals during emergencies. The researchers utilized ISO/IEC 25010 to
evaluate software quality with 243 respondents.

Another related study is \emph{READI: A Modern Way of Saving Lives Using
	Mobile Application: Real-Time Emergency Alerts and Disaster Information}
\parencite{deguzman2019}. The researchers developed a mobile application providing
real-time emergency alerts and awareness materials to the public, emphasizing
the importance of mobile technology in improving communication during
emergencies.

While both studies successfully provided a technological basis for disaster
management, they focused strictly on the manual emergency response and
communication aspects. Neither study attempted to implement natural language
processing pipelines, such as contextual Joint Matrices or generative AI Large
Language Models (LLMs), to automatically analyze, cluster, and score sentiments
on the unstructured text data inherent to disaster communications.

\section{Synthesis of Literature and Studies, Research Gaps, and Justification}

Literature at the international, national, and institutional levels
demonstrates the effective use of probabilistic modeling, specifically Latent
Dirichlet Allocation (LDA) and contextual transformers, in processing
unstructured data into highly interpretable thematic clusters.

Although existing studies have demonstrated significant advancements, critical
computational gaps remain. The institutional studies reviewed rely primarily on
sensor data, manual user inputs, and GPS hardware, leaving them unequipped to
automatically parse high-velocity, multilingual disaster reports. Conversely,
international studies successfully integrate Natural Language Processing (NLP)
and contextual embeddings, but these methodologies have yet to be rigorously
benchmarked for localized Philippine deployment. Local Disaster Risk Reduction
and Management Offices (LDRRMOs) currently lack the computational
infrastructure to process mixed-syntax Taglish text at scale.

The present study seeks to bridge this gap through a novel comparative
methodology. Instead of relying on classical classifiers or manual GPS
tracking, the proposed framework introduces a Joint Matrix architecture
integrating LDA topic probabilities with \xlmr{} contextual embeddings to
preserve interpretability while mapping urgent distress signals. By benchmarking
this hybrid system against a state-of-the-art open-weight dense large language
model (Qwen3-1.7B-Instruct) operating under few-shot conditions, this research
provides the necessary empirical validation to integrate advanced, automated
semantic and sentiment modeling of natural disaster-related social media posts.